% 18-research.tex: what is held, what is asked, and what is not the repository's to do.
\section{Future research and engineering, under a rule that forbids most of it}
\label{sec:research}

Version 2's version of this section was a list of things to build. This one cannot be, because on
13 September 2026 the repository adopted an operating contract that changes what \emph{future work}
means. A product change now needs exactly one qualifying reason: a named external implementation
requires it, a named non-author operator or holder requires it, a finding originating outside the
repository requires it, or an executable counterexample shows behaviour the repository already
promises to be false. Cleanliness, completeness, elegance and a desirable new guarantee do not
qualify, and the contract says not to invent a fifth reason. Everything else goes to the lab, whose
job is to falsify the differentiating claims, not to expand the architecture. So the items below are
sorted by who would have to act, and everything is marked \sFut{} unless stated otherwise.

\subsection{Gates that are not the repository's to pass}

The contract's ninety-day objective names them, and the scoreboard records each as filled or blank:
an operator who is not the author reaching a verified result without help, which is the one thing
that separates the release candidate from \code{1.0.0}; a relying party that is not the author; a
second wallet, and a wallet on the native rather than the OpenID4VP path; a certification submission
that would put the four modules in \code{REVIEW} in front of a Foundation reviewer; an independent
security review, penetration test or cryptographic audit; a pilot with consenting users; and
institutionally independent witnesses. The contract also names the other end. If, at 180 days, the
external rows remain effectively empty while internal versions, invariants and architecture keep
growing, the core is to be frozen, the national-scale architecture stops expanding, the verifier is
kept only if external use exists, and the project is archived as a reference implementation. This
paper's subject may be archived by its author's own rule, and the paper records the rule rather than
arguing with it. \sExt

\subsection{Questions the lab still owes an answer}

Each of these is a measurement, not a feature, and the lab's discipline is to publish the threat
model, the adversary, the harness, the result and the limitations together, with a positive control
before any null result counts.

\begin{description}
  \item[The rest of the linkability threat model.] The measured result covers stable fields
    only. Presentation size, timing, issuer metadata narrowing the anonymity set, status artifacts
    stapled to a presentation, and transcript structure are unmeasured, and the one channel the
    adversary can read, transcript length, is known to be live in any deployment where a
    presentation's size varies per holder.
  \item[The interoperable path's nine handles.] Batch issuance, many credentials with one holder key
    each, is the ecosystem's mitigation and is not implemented; whether it would be adopted is an
    outside party's question, and the cost of not having it is now a number.
  \item[The migration path's other half.] Rollback, and the mixed window during a partial
    migration, are unmeasured. A verifier-readable algorithm status, in the registry or the trust
    list, would be a new product guarantee, so the lab recorded the gap and did not build it.
  \item[Duress against the informed coercer.] A holder-side panic signal, a deniable enrolment that
    makes the absence of a duress code unprovable, and a purge path for the duress record are each a
    mechanism change, and the assessment proposes none of them. What it asks for first is a
    measurement the repository cannot make: what a coercer standing beside a holder actually
    observes.
  \item[The anonymity-set floor.] Reporting the crowd size beside a \code{bounded} verdict, or
    refusing to close an epoch below a floor, would change what the system promises; both are
    recorded for a deploying organisation to decide.
\end{description}

\subsection{Presentation technologies that would remove the correlator, held}

Version 2 listed five families. One, the pairwise identifier, is built, with the bound
Section~\ref{sec:privacy} states. The other four are held rather than scheduled, and the reason each
is not built is the contract's, not a technical one.

\begin{description}
  \item[Blinded or derived presentations.] The holder proves possession of a signature over
    attributes without revealing the signature itself. The classical constructions rely on
    pairing-based assumptions that a quantum computer breaks; post-quantum equivalents are an open
    research area, and the repository's posture would have to be preserved.
  \item[One-time presentation tokens.] Single-use artifacts minted per presentation, at the cost of
    issuer contact or a pre-fetched batch, which trades against offline availability and issuer
    non-observation exactly as the status assertion does.
  \item[Anonymous credentials with selective disclosure.] Revealing that an attribute satisfies a
    predicate without revealing the attribute. The benchmark against a mature system measured the
    gap precisely: Polaris as an issuer has no attribute credential to apply a predicate to, and as a
    verifier it accepts a pre-computed boolean the issuer had to anticipate.
  \item[A native anonymous predicate.] The benchmark's first scenario is satisfied by a disclosed
    attribute, not by an anonymous predicate; that is on the scoreboard as a known limitation, which
    moves only when an outside party trips over it.
\end{description}

\subsection{Metis: learning the system, never the people}

The assessment the roadmap requires before any code would have to falsify the central hypothesis
(that a model can infer properties of the system without inferring the rights of people) against a
synthetic-nation ground truth, compare every model to a deterministic rule, and design the
constitutional invariant that would forbid the prohibited uses structurally, in the manner of
Athena's five checks. Under the contract the assessment itself waits on an outside party needing
what it would produce; nothing about Metis is scheduled, and this paper keeps the brief because the
prohibition in it binds any future cognitive layer whether or not Metis is ever built.

\subsection{Myrmex: stigmergic debugging as an outside observer}

Version 2 proposed, under the provisional name Myrmex, a debugging system that would live outside
Polaris and attack it independently, inspired by how an ant colony coordinates without a central
plan: specialised agents leaving machine-readable evidence on a shared landscape, confidence drawn
only from reproductions and never from agreement between models, and every reproducible failure
entering the loop of Section~\ref{sec:senses} as a test and an invariant. The proposal stands, and
the caution with it: an earlier apparatus that lived inside the tree and asserted claims about itself
was deleted as scaffolding, so Myrmex would be held to the opposite discipline, outside and
adversarial, or not built.

\figfile{stigmergy}{The proposed colony. Small specialised agents traverse the code and the
architecture and leave machine-readable pheromones on a persistent evidence landscape: suspicion, a
trust-boundary anomaly, an unusual state transition, a cryptographic assumption, a missing edge, an
inconsistent claim, an area repeatedly associated with failure. Pheromones strengthen under
independent evidence and decay unsupported; agents follow reinforced trails. An orchestrator allocates
work and aggregates evidence and decides nothing. Confidence comes only from external evidence, never
from agreement between models. What survives reproduction enters Polaris's own loop as a test and an
invariant.}{fig:stigmergy}

What Version 3 can add is that the thing Myrmex was for has been demonstrated without it. The
mutation drills are the colony's cheapest members, built inside the tree: they do not reason, they
delete and invert and watch for red, and they found more than any review had. And the observer that
found what no drill could, a certificate every internal instrument accepted, was not a colony of
agents but one wallet written by strangers. The proposal is kept because a stranger's wallet
exercises one path, and the rings inside it are still watched only by their author's instruments.

\subsection{Engineering items held under the contract}

A hash-based signer behind the same custody interface; the chain-backed ledger drivers whose
interfaces are declared; native wallet clients and a browser bridge; status distribution as a
deployment rather than a specification; a manufactured card; and the external-implementation path
for the native protocol. Each is a roadmap row that the contract has moved from \emph{next} to
\emph{when somebody outside needs it}, and the repository's own history is the argument for the
move: three hundred and ninety-seven versions were spent building rings that no one outside had
asked for, and the first outside party to arrive found a defect in an afternoon.
